\section{Analysis}

\subsection{Why Does HGRN2Layer Work?}

\subsubsection{Element-Wise State is Efficient}

The state $s_t \in \mathbb{R}^d$ is a vector, not a matrix. Each dimension 
updates independently. This gives $\mathcal{O}(L d)$ complexity instead of 
$\mathcal{O}(L H d^2)$ for matrix-state recurrences (DeltaNet).

\subsubsection{Bounded Forget Gate}

The clamp $\log f \in [-3, -10^{-4}]$ keeps $f \in [0.05, 1]$. Without 
clamping:

\begin{itemize}
\item Too small $f$ $\to$ state collapses to zero (no memory).
\item Too large $f$ $\to$ state saturates (unstable).

\end{itemize}
\textbf{Ablation:} Removing the clamp drops accuracy by {{CLAMP\_DROP}}\%.

\subsubsection{Input Modulation}

Multiplying the input by $(1 - f_t)$ ties the write magnitude to the forget 
rate. This prevents the state from drifting away from the input.

\textbf{Ablation:} Removing the modulation drops accuracy by {{MODULATION\_DROP}}\%.

\subsubsection{Bidirectional for Encoders}

DistilBERT and RoBERTa use bidirectional scan; GPT-2 uses causal. The 
bidirectional scan improves accuracy by {{BIDI\_GAIN}}\% on DistilBERT.

\subsection{Where Does the Slowdown Come From?}

\subsubsection{Python Loop Over Chunks}

The scan is chunked but still uses a Python loop over chunks:

\begin{lstlisting}
for start in range(0, L, chunk_size):
    ...
\end{lstlisting}

For L = 128 and chunk = 32, that is 4 iterations. Across 6 layers and 109 
batches, that is ~2600 Python iterations.

\subsubsection{Intermediate Tensors}

Each chunk materializes \texttt{A}, \texttt{A\_diff}, \texttt{decay\_weights}, \texttt{intra} tensors, all 
in FP32. These accumulate in GPU memory.

\subsubsection{Why Long Sequences Are Faster}

At L = 4096, the baseline's $\mathcal{O}(L^2)$ attention dominates the cost. 
Our scan is $\mathcal{O}(L)$ -- it scales linearly, so the crossover happens.

\subsection{Per-Layer Analysis}

\textbf{Table 8.1: Forget gate statistics per layer.}

\begin{table}[H]
\centering
\small
\begin{tabular}{lllll}
\toprule
\textbf{Layer} & \textbf{Mean f} & \textbf{Std f} & \textbf{Min f} & \textbf{Max f} \\
\midrule
0 & {{D\_L0\_F\_MEAN}} & {{D\_L0\_F\_STD}} & {{D\_L0\_F\_MIN}} & {{D\_L0\_F\_MAX}} \\
3 & {{D\_L3\_F\_MEAN}} & {{D\_L3\_F\_STD}} & {{D\_L3\_F\_MIN}} & {{D\_L3\_F\_MAX}} \\
5 & {{D\_L5\_F\_MEAN}} & {{D\_L5\_F\_STD}} & {{D\_L5\_F\_MIN}} & {{D\_L5\_F\_MAX}} \\
\bottomrule
\end{tabular}
\end{table}

\subsection{Comparison to Attention Baselines}

\textbf{Table 8.2: Accuracy comparison.}

\begin{table}[H]
\centering
\small
\begin{tabular}{lllll}
\toprule
\textbf{Method} & \textbf{DistilBERT} & \textbf{RoBERTa} & \textbf{GPT-2} & \textbf{Speed @ L=128} \\
\midrule
full & 90.71\% & 78.56\% & 79.36\% & 1$\times$ \\
Linear Transformer & {{LT\_D}} & {{LT\_R}} & {{LT\_G}} & {{LT\_SP}}$\times$ \\
Performer & {{PERF\_D}} & {{PERF\_R}} & {{PERF\_G}} & {{PERF\_SP}}$\times$ \\
\textbf{HGRN2Layer} & \textbf{{{D\_ACC}}} & \textbf{{{R\_ACC}}} & \textbf{{{G\_ACC}}} & \textbf{{{D\_SP\_128}}$\times$} \\
\bottomrule
\end{tabular}
\end{table}

\subsection{Failure Analysis}

\subsubsection{RoBERTa Degradation}

RoBERTa shows a larger drop than DistilBERT. Hypotheses:

\begin{enumerate}
\item \textbf{Nested handling.} \texttt{is\_nested\_encoder=True} sets
  \texttt{layer.attention.output.dense = nn.Identity()}. If this is wrong,
  output projection is skipped.
\item \textbf{12 layers.} More layers to accumulate error.

\end{enumerate}
\textbf{Test.} {{ROBERTA\_TEST}}

\subsubsection{GPT-2 Causal}

GPT-2 is causal, so only the forward scan is used. This is stricter than 
bidirectional.

\subsection{Summary of Insights}

\begin{enumerate}
\item \textbf{Element-wise recurrence works} at L $\ge$ {{BREAKEVEN}}.
\item \textbf{Bounded forget gate is essential} ({{CLAMP\_DROP}}\% without).
\item \textbf{Input modulation prevents drift} ({{MODULATION\_DROP}}\% without).
\item \textbf{Bidirectional helps encoders} by {{BIDI\_GAIN}}\%.
\item \textbf{Python chunk loop limits short-sequence speed.}
\end{enumerate}
\begin{lstlisting}

---
\end{lstlisting}